%%%%%%%%%%%%%%%%%%%%%%%%%%%%%%%%%%%%%%%%%%%%%%%%%%%%%%%%%%%%
\section*{Appendix B: Custom Verifier Implementation}
\label{appendix:verifiers}

\subsection*{Verification Pipeline Architecture}

{\tiny
The PEP verification pipeline composes five stages:

\begin{verbatim}
Invocation Arrives at PEP
         ↓
┌──────────────────────────────────────────────────┐
│  Stage 1-3: DETERMINISTIC (100% recall, 0% FP)  │
│  • Signature verification (Lemma 1)             │
│  • Policy binding verification                   │
│  • MAPL policy evaluation (Lemma 4)             │
│                                                   │
│  If this fails -> REJECT (cryptographically)    │
└─────────────────┬────────────────────────────────┘
                  ↓ PASS (authenticated + authorized)
┌──────────────────────────────────────────────────┐
│  Stage 4: CUSTOM VERIFIERS (may have FPs)       │
│  • MemoryIntegrityVerifier                       │
│  • WorkflowIntegrityVerifier                     │
│  • ToolAuthorizationVerifier                     │
│  • StorageIntegrityVerifier                      │
│  • + 18 more (Phases 2-4)                       │
│                                                   │
│  Operates on VERIFIED context                    │
│  Can only ADD restrictions, never remove         │
│                                                   │
│  If this fails -> REJECT (policy violation)     │
└─────────────────┬────────────────────────────────┘
                  ↓ PASS
┌──────────────────────────────────────────────────┐
│  Execute Operation                               │
└─────────────────┬────────────────────────────────┘
                  ↓
┌──────────────────────────────────────────────────┐
│  Stage 5: POST-PROCESSING VERIFIERS              │
│  • PIIMaskingVerifier                            │
│  • Secret removal                                │
│  • Output validation                             │
└──────────────────────────────────────────────────┘
\end{verbatim}

Even if custom verifiers are compromised, Stages 1-3 enforce cryptographic integrity and policy constraints.
}

\subsection*{Production Verifiers}

Four production verifiers address OWASP Top 10 threats through domain-specific checks on cryptographically verified context:

{\tiny
\noindent\begin{minipage}[t]{0.48\columnwidth}
\textbf{MemoryIntegrityVerifier}\\
\textit{Prevents memory poisoning (OWASP LLM01)}\\
• Rate limiting (10 updates/min)\\
• Protected fields (goals, system\_prompt)\\
• Pattern detection (suspicious updates)\\
• Cryptographic integrity checks
\end{minipage}
\hfill
\begin{minipage}[t]{0.48\columnwidth}
\textbf{WorkflowIntegrityVerifier}\\
\textit{Prevents workflow hijacking (OWASP LLM04)}\\
• Sequence enforcement (prerequisites)\\
• State tracking (prevent re-execution)\\
• Attestation chains (crypto proof)\\
• Duration limits (max 3600s)
\end{minipage}

\vspace{0.5em}

\noindent\begin{minipage}[t]{0.48\columnwidth}
\textbf{ToolAuthorizationVerifier}\\
\textit{Prevents tool misuse (OWASP LLM06)}\\
• RBAC (role-to-tool mapping)\\
• Dangerous patterns (write→execute)\\
• Usage rate limiting (30 tools/min)\\
• Privilege escalation prevention
\end{minipage}
\hfill
\begin{minipage}[t]{0.48\columnwidth}
\textbf{StorageIntegrityVerifier}\\
\textit{Prevents data exfiltration (OWASP LLM02)}\\
• Path traversal prevention\\
• Data classification (secret/sensitive)\\
• Encryption enforcement\\
• Exfiltration detection (bulk reads)
\end{minipage}
}
